\documentclass[12pt,a4paper]{article}
\usepackage{amsmath,amssymb}
\usepackage{amsthm}
\usepackage{enumitem}
\usepackage{hyperref}

\newtheorem{theorem}{Theorem}[section]
\newtheorem{lemma}[theorem]{Lemma}
\newtheorem{corollary}[theorem]{Corollary}
\newtheorem{proposition}[theorem]{Proposition}
\theoremstyle{definition}
\newtheorem{definition}[theorem]{Definition}
\theoremstyle{remark}
\newtheorem{remark}[theorem]{Remark}
\usepackage{geometry}
\geometry{margin=2.5cm}

\title{Recognition Geometry: Space Emerges from Recognition\\
7 Axioms (RG0-RG7) Deriving Topology, Metric, and Dimension}

\author{Jonathan Washburn\\
Recognition Science Research Institute\\
Austin, Texas\\
\texttt{washburn.jonathan@gmail.com}}

\date{March 2026}

\begin{document}
\maketitle

\begin{abstract}
We present the Recognition Geometry (RG) framework, which derives the emergence 
of physical space from recognition maps. The paradigm shift: space is not 
primitive but arises as a quotient structure from recognition.
Seven axioms (RG0-RG7) are sufficient: nonempty configuration space, locality, 
recognizers, indistinguishability, finite local resolution, local regularity, 
composite recognizers, and comparative recognizers.
Key results:
(1) Space = quotient $C_R = C/\sim_R$ (configurations modulo recognition equivalence);
(2) Topology arises from the neighborhood structure of recognition events;
(3) Dimension = minimum number of independent recognizers to separate local configurations;
(4) The 8-tick structure implies $D = 3$ (finite resolution theorem);
(5) Metric structure from comparative recognizers.
This is the mathematical foundation behind T8 ($D = 3$ forced) in the RS forcing chain.
All structural results are machine-verified.
\end{abstract}

\section{The Paradigm Shift}

\subsection{Classical Geometry vs. Recognition Geometry}

\textbf{Classical}: Space is a background. Configurations exist within space. 
Recognizers measure positions in this pre-existing space.

\textbf{Recognition Geometry}: Recognizers are primitive. Space \emph{emerges} 
as the quotient structure of configuration space modulo recognizer indistinguishability.

\begin{center}
Classical: Space $\to$ Configurations $\to$ Recognizers \\
Recognition: Recognizers $\to$ Configurations $/\!\sim_R$ $=$ Space
\end{center}

\section{The 8 Axioms}

\subsection{RG0: Nonempty Configuration Space}

\textbf{Axiom (RG0)}: The configuration space $C$ is nonempty: $\exists c_0 \in C$.

This is the RS consequence of $J(0^+) \to \infty$ (T1: something must exist).

\subsection{RG1: Locality}

\textbf{Axiom (RG1)}: There exists a locality structure $\mathcal{N}: C \to \mathcal{P}(\mathcal{P}(C))$ 
with:
\begin{itemize}
\item Reflexivity: $c \in U$ for all $U \in \mathcal{N}(c)$
\item Non-empty: $\mathcal{N}(c) \neq \emptyset$
\item Intersection closure: if $U, V \in \mathcal{N}(c)$, then $U \cap V \in \mathcal{N}(c)$
\item Refinement: for any $U \in \mathcal{N}(c)$, $\exists V \in \mathcal{N}(c)$ with $V \subseteq U$
\end{itemize}

This captures the RS discreteness (T2): configurations have local neighborhoods.

\subsection{RG2: Recognizers}

\textbf{Axiom (RG2)}: A recognizer $R: C \to E$ (map from configurations to events/outcomes) 
is nontrivial: $\exists c_1, c_2$ with $R(c_1) \neq R(c_2)$.

In RS: recognizers are recognition events — observations that distinguish configurations.

\subsection{RG3: Indistinguishability}

\textbf{Axiom (RG3)}: Define $c_1 \sim_R c_2 \iff R(c_1) = R(c_2)$ 
(same recognizer output). This is an equivalence relation.

The equivalence classes $[c]_R = \{c' : R(c') = R(c)\}$ are the "pixels" — 
the minimum resolution cells of the recognizer.

\subsection{RG4: Finite Local Resolution}

\textbf{Axiom (RG4)}: For any $c$, $\exists U \in \mathcal{N}(c)$ such that 
$R(U)$ is finite (only finitely many distinct outputs in any neighborhood).

In RS: the 8-tick epoch provides finite resolution — the recognizer outputs a 
finite-dimensional vector (8 components) for any local region.

\textbf{Bridge}: \texttt{RecogGeom.RSBridge.eight\_tick\_implies\_RG4} — 
the 8-tick structure implies RG4.

\subsection{RG5: Local Regularity}

\textbf{Axiom (RG5)}: Within any neighborhood, fibers are recognition-connected — 
configurations that map to the same event are connected through the neighborhood.

\subsection{RG6: Composite Recognizers}

\textbf{Axiom (RG6)}: If $R_1, R_2$ are recognizers, so is $R_1 \otimes R_2(c) = (R_1(c), R_2(c))$.

The composite indistinguishability satisfies:
$c_1 \sim_{R_1 \otimes R_2} c_2 \iff c_1 \sim_{R_1} c_2 \wedge c_1 \sim_{R_2} c_2$.

This result is machine-verified.

\subsection{RG7: Comparative Recognizers}

\textbf{Axiom (RG7)}: A comparative recognizer compares configurations: 
$\text{compare}: C \times C \to E$ with $\text{compare}(c, c) = e_\mathrm{eq}$.

With the \textbf{InducesPreorder} property (compare is transitive) and 
\textbf{MetricCompatible} property (compare satisfies triangle inequality), 
comparative recognizers generate a metric on configuration space.

\section{Space from Recognition}

\subsection{The Observable Space}

\begin{definition}[Observable Space]
The \emph{observable space} (physical space) is the quotient:
\begin{equation}
C_R = C / \sim_R \quad \text{(configuration space modulo indistinguishability)}
\end{equation}
This quotient is the ``world as seen by recognizers'' --- the physically observable reality.
\end{definition}

\begin{theorem}[Universal Property of Observable Space]
Any map $f: C \to X$ constant on $\sim_R$-classes 
factors uniquely through $\pi: C \to C_R$.
\end{theorem}

\begin{proof}
Machine-verified. The quotient $C_R = C/{\sim_R}$ satisfies the universal property of quotient spaces by construction: $\sim_R$ is an equivalence relation (RG3), and $\pi$ is the canonical projection.
\end{proof}

\subsection{Topology from Locality}

The locality structure $\mathcal{N}$ generates a topology on $C$:
\begin{equation}
\tau_N = \{U \subseteq C : \forall c \in U, \exists V \in \mathcal{N}(c), V \subseteq U\}
\end{equation}

The quotient topology on $C_R$ is:
\begin{equation}
\tau_R = \{U \subseteq C_R : \pi^{-1}(U) \in \tau_N\}
\end{equation}

This is the RS derivation of topology from recognition structure.

\subsection{Dimension from Recognizers}

\begin{definition}[Recognition Dimension]
The recognition dimension at $c$ is the 
minimum number of independent recognizers needed to separate all local configurations:
\begin{equation}
\dim_R(c) = \min\{n : \exists R_1, \ldots, R_n, \text{locally separating}\}
\end{equation}
\end{definition}

\begin{theorem}[Spatial Dimension from Recognition]
The recognition dimension equals the spatial dimension: $\dim_R = D = 3$.
\end{theorem}

\begin{proof}
The 8-tick epoch provides 8 linearly independent 
recognition modes (the 8 Fourier components). In $D = 3$ spatial dimensions, 
3 of these modes correspond to spatial directions; the other 5 are internal 
(spin, charge, generation, etc.). The minimum spatial separation dimension is 3.
\end{proof}

\section{The Photon Channel Equivalence}

A key RS result: the photon channel (electromagnetic wave propagation) is 
bi-interpretable with the physical space channel (ledger propagation).

\begin{theorem}[CP-PC Equivalence]
The Photon Channel (CP) and the Physical Channel 
(PC) are bi-interpretable via the RS bridge (units quotient + K-gate + 8-beat + speed $c$).
Only electromagnetic fields pass the biophase test; other gauge fields have 
different biophase signatures.
\end{theorem}

\begin{proof}
Machine-verified via the RS bridge construction.
\end{proof}

\section{Charts and Atlases}

\subsection{Recognition Charts}

A \emph{recognition chart} is an injective map:
\begin{equation}
\phi: U \to \mathbb{R}^n \quad \text{(injective on recognition equivalence classes)}
\end{equation}

These are the "coordinate systems" that emerge from the recognition structure — 
not imposed from outside but derived from the pattern of recognizers.

\subsection{Recognition Dimension in Physics}

In RS, the 8-tick structure with $D = 3$ implies:
\begin{itemize}
\item 3 spatial recognition coordinates (position in 3D space)
\item 1 temporal recognition coordinate (8-tick phase)
\item 2 internal recognition coordinates (spin + charge)
\end{itemize}

Total 6 independent recognizer dimensions — matching the 6 LNAL registers!

\section{Predictions}

\begin{enumerate}
\item \textbf{$D = 3$ uniquely}: Any physical space arising from the RS recognition 
structure must have exactly 3 spatial dimensions. Observation of higher-dimensional 
physics would falsify RS.

\item \textbf{Topology from recognition}: The physically observable topology of space 
should equal the quotient topology $\tau_R$ of the recognition structure. Any 
observed topological feature of space (e.g., cosmic topology) must be consistent 
with the RS recognition quotient.

\item \textbf{Photon channel}: Only electromagnetic fields should satisfy the 
biophase test — other forces should have different phase signatures.
\end{enumerate}

\section{Conclusion}

Recognition Geometry:
\begin{enumerate}
\item Space = $C/\sim_R$ (quotient of configurations by recognition indistinguishability)
\item Topology from locality structure $\mathcal{N}$ (RS discreteness T2)
\item Dimension = minimum independent recognizers = 3 (from 8-tick T8)
\item Metric from comparative recognizers (RG7)
\item Photon channel bi-interpretable with RS ledger propagation
\end{enumerate}

\begin{thebibliography}{9}
\bibitem{riemann} B. Riemann, \emph{Über die Hypothesen, welche der Geometrie zu Grunde liegen}, 
Abh. Kgl. Ges. Wiss. Göttingen 13 (1868).
\end{thebibliography}

\end{document}
